% Relatório para o autor entender o Hannan-Rissanen (30/10/2026, pedido do autor).
% Material de estudo, não é parte do TG. Todo número vem de
% dados/resultados/hannan_rissanen_2026-10-30/ (ferramentas/scripts/hannan_rissanen.sh)
% ou de dados/resultados/robo_em_acao_2026-10-30/ (traço da refação de 05/10/2026).
\documentclass[11pt,a4paper]{article}
\usepackage[T1]{fontenc}
\usepackage[brazilian]{babel}
\usepackage{amsmath,amssymb}
\usepackage{mathptmx}
\usepackage[a4paper,margin=2.4cm]{geometry}
\usepackage{booktabs,array}
\usepackage[font=small,labelfont=bf]{caption}
\usepackage{enumitem}
\usepackage[most]{tcolorbox}
\usepackage{pgfplots}
\usepgfplotslibrary{groupplots,dateplot,statistics}
\usetikzlibrary{arrows.meta,positioning,babel,shapes.misc}
\pgfplotsset{compat=1.18}
\usepackage[hidelinks]{hyperref}
\setlength{\parskip}{0.45em}
\setlength{\parindent}{0pt}

\definecolor{cmod}{RGB}{0,90,181}
\definecolor{cdes}{RGB}{220,50,32}
\definecolor{cdad}{RGB}{120,120,120}
\definecolor{cari}{RGB}{230,140,0}
\definecolor{cok}{RGB}{0,140,70}

\newcommand{\arq}[1]{../../../dados/resultados/hannan_rissanen_2026-10-30/#1}
\newcommand{\arqr}[1]{../../../dados/resultados/robo_em_acao_2026-10-30/#1}
\input{\arq{valores.tex}}
\input{\arq{real_matrizes.tex}}

\pgfplotsset{base/.style={table/col sep=comma, grid=major, grid style={gray!20},
    tick label style={font=\small}, label style={font=\small},
    title style={font=\small\bfseries, align=center},
    every axis plot/.append style={line width=0.9pt},
    /pgf/number format/use comma, /pgf/number format/1000 sep={.}}}

\tikzset{fluxo/.style={draw=#1, fill=#1!7, rounded corners=3pt, text width=3.0cm, align=center, minimum height=1.35cm, font=\footnotesize}}
\newtcolorbox{ideia}[1][]{colback=cmod!5, colframe=cmod!70, fonttitle=\bfseries, title={#1}, arc=2pt, boxrule=0.6pt}
\newtcolorbox{cuidado}[1][]{colback=cdes!4, colframe=cdes!70, fonttitle=\bfseries, title={#1}, arc=2pt, boxrule=0.6pt}
\newtcolorbox{banca}{colback=black!3, colframe=black!50, fonttitle=\bfseries, title=Se a banca perguntar, arc=2pt, boxrule=0.6pt}

\title{\textbf{O método de Hannan-Rissanen no seu ARIMA}\\[2pt]\large como ele funciona, passo a passo, e o que ele fez no robô}
\author{Material de estudo para os autores do TG (não faz parte do texto)}
\date{30 de outubro de 2026}

\begin{document}
\maketitle

\begin{ideia}[Em uma página]
\textbf{O que é.} Uma forma de estimar os coeficientes de um ARIMA usando só duas regressões de mínimos quadrados
(MQO), sem otimização iterativa.

\textbf{O problema que resolve.} O termo de médias móveis (MA) do ARIMA usa o \emph{erro} de ontem, e o erro não
aparece nos dados. Sem ele não dá para montar a regressão.

\textbf{A solução, em duas frases.} Primeiro estima-se um modelo autorregressivo longo, cujos resíduos servem de
estimativa dos erros (estágio 1). Depois esses erros estimados entram como mais uma coluna numa segunda
regressão, que dá os coeficientes do modelo (estágio 2).

\textbf{No seu robô.} Refazendo fora do robô a estimação de dois pregões reais, os coeficientes e a previsão saíram
\emph{idênticos} aos que o robô registrou (diferença zero). No traço de jul/2019 a jun/2021, os dois
estágios rodaram em \vHTracoDois{} dos \vHTracoN{} pregões ($\vHTracoPct\%$). Nos outros, o modelo escolhido não
tinha termo MA e o método se reduz a um MQO comum. Mas na escolha das ordens, a cada 21 pregões, 20 dos 25
modelos candidatos passam pelos dois estágios.

\textbf{O que se descobriu.} Em séries simuladas o método acerta perto do valor verdadeiro, com um pouco mais de
erro que a máxima verossimilhança. Para um $\theta$ do tamanho do que o robô estimou no pregão detalhado, o erro de
estimação é quase do tamanho do próprio coeficiente.
\end{ideia}

\tableofcontents

%=====================================================================
\section{O problema: o erro de ontem não está nos dados}

O ARIMA do robô, depois de diferenciar o log-preço ($d = 1$), descreve o log-retorno $z_t$ de cada dia assim
(é a convenção do código, \texttt{Arima.mqh}):
\[
  z_t \;=\; c \;+\; \underbrace{\phi_1 z_{t-1} + \dots + \phi_p z_{t-p}}_{\text{parte AR: retornos passados}}
  \;+\; \underbrace{\theta_1 \varepsilon_{t-1} + \dots + \theta_q \varepsilon_{t-q}}_{\text{parte MA: erros passados}}
  \;+\; \varepsilon_t .
\]
O $\varepsilon_t$ é o erro do dia: a parte do retorno que o modelo não consegue prever.

Se o modelo só tivesse a parte AR, a estimação seria uma regressão comum. Cada linha da tabela seria um dia, a
coluna da esquerda seria $z_t$ e as da direita seriam $z_{t-1}, z_{t-2}, \dots$, todos números conhecidos. O MQO
escolhe os $\phi$ que deixam a soma dos erros ao quadrado a menor possível, numa conta fechada.

Com a parte MA aparece o problema: para montar a coluna $\varepsilon_{t-1}$ é preciso saber o erro de ontem, e o erro
depende dos próprios coeficientes que se quer estimar. É uma volta em círculo:

\begin{center}
\begin{tikzpicture}[caixa/.style={draw, rounded corners=3pt, fill=#1!8, text width=4.2cm, align=center, minimum height=1cm, font=\small},
  seta/.style={-{Stealth[length=6pt]}, line width=0.9pt}]
\node[caixa=cdes] (a) {para estimar $\theta$\\preciso da coluna $\varepsilon_{t-1}$};
\node[caixa=cdes, right=2.2cm of a] (b) {para calcular $\varepsilon_{t-1}$\\preciso de $\theta$};
\draw[seta] (a.north east) to[bend left=25] node[above, font=\small] {?} (b.north west);
\draw[seta] (b.south west) to[bend left=25] node[below, font=\small] {?} (a.south east);
\end{tikzpicture}
\end{center}

Há duas saídas. A primeira é a \textbf{máxima verossimilhança}: chutar coeficientes, calcular os erros, medir quão
plausíveis eles são e repetir até não melhorar mais. Funciona, mas é iterativa: precisa de ponto de partida, de
critério de parada e pode não convergir. A segunda é o \textbf{Hannan-Rissanen}, que quebra o círculo estimando os
erros \emph{antes}, por um caminho que não depende de $\theta$.

\begin{cuidado}[E se simplesmente ignorar a parte MA?]
Numa simulação com $\phi = \vHSintPhi$ e $\theta = \vHSintTheta$ verdadeiros (seção~\ref{sec:montecarlo}), um AR(1) por MQO
que ignora o MA estimou, em média, $\phi = \vHMcArmaIngenuo$ em vez de $\vHSintPhi$. O efeito do MA não some: ele vai
parar dentro do coeficiente AR e o distorce.
\end{cuidado}

%=====================================================================
\section{A ideia que destrava tudo: um ARMA é um AR comprido}

Se a parte MA for \emph{invertível} ($|\theta| < 1$ no caso de um termo), dá para reescrever o erro de hoje como
uma soma de retornos passados. Para um ARMA(1,1), substituindo $\varepsilon_{t-1}$ pela própria equação de ontem,
e depois $\varepsilon_{t-2}$, e assim por diante, chega-se a
\[
  z_t = \pi_1 z_{t-1} + \pi_2 z_{t-2} + \pi_3 z_{t-3} + \dots + \varepsilon_t,
  \qquad \pi_j = (\phi + \theta)\,(-\theta)^{\,j-1}.
\]
Ou seja, um ARMA(1,1) é um AR com infinitos termos, cujos pesos $\pi_j$ encolhem a cada defasagem. Um AR
\emph{comprido, mas finito} chega bem perto dele. E um AR se estima por MQO, sem precisar de $\varepsilon$.

A Figura~\ref{fig:pesos} mostra isso numa série simulada com $\phi = \vHSintPhi$ e $\theta = \vHSintTheta$, do mesmo tamanho
da janela do robô ($\vHSintN$ retornos). As barras azuis são os pesos verdadeiros; os pontos laranja, o que o AR
longo do estágio 1 estimou a partir dos dados.

\begin{figure}[htbp]\centering
\begin{tikzpicture}
\begin{axis}[base, width=14cm, height=5.6cm, xmin=0, xmax=30, ymin=-0.6, ymax=1.1, xlabel={Defasagem $j$},
  ylabel={Peso $\pi_j$}, /pgf/number format/fixed,
  legend style={at={(0.98,0.97)}, anchor=north east, font=\footnotesize}]
\addplot[ybar, bar width=5pt, fill=cmod!70, draw=none] table[x=j, y=verdadeiro]{\arq{sint_pesos.csv}};
\addlegendentry{peso verdadeiro $(\phi+\theta)(-\theta)^{j-1}$}
\addplot[only marks, mark=*, mark size=1.8pt, cari] table[x=j, y=estimado]{\arq{sint_pesos.csv}};
\addlegendentry{estimado pelo AR($\vHSintM$) do estágio 1}
\end{axis}
\end{tikzpicture}
\caption{Série simulada: o AR longo recupera os primeiros pesos, que são os que importam. Depois da quinta
defasagem os pesos verdadeiros são quase zero, e o que o AR estima ali é ruído.}
\label{fig:pesos}
\end{figure}

%=====================================================================
\section{Estágio 1: o AR longo que estima os erros}

\subsection{Quantas defasagens}
O código escolhe a ordem $m$ do AR longo pela regra
\[
  m = \lfloor (\ln N)^2 \rfloor, \qquad \text{limitada por baixo a } p+q+1 \text{ e por cima a } N/5,
\]
em que $N$ é o número de retornos. A janela do robô tem 240 fechamentos e portanto $N = \vHRealN$ retornos.
Com isso, $m = \vHRealM$. A regra vem da literatura: Gómez e Maravall (1998,
p.~28) usam $\lfloor \log^2 N \rfloor$, apoiados em Hannan e Rissanen (1982, p.~88), que pedem que $m$ cresça com o
tamanho da série, mas devagar.

\subsection{A regressão}
Cada linha é um dia $t$ a partir do $(m+1)$-ésimo, porque os primeiros $m$ dias não têm $m$ dias anteriores.
São $N - m = \vHRealObsUm$ linhas e $m + 1 = \the\numexpr\pHRealM+1\relax$ colunas (a constante e os $\vHRealM$ retornos anteriores). A tabela
mostra as primeiras linhas no pregão real que será detalhado na seção~\ref{sec:real} (retornos em \%):

\begin{center}\small
\begin{tabular}{l r | c r r r c r}
\toprule
dia $t$ & $z_t$ & 1 & $z_{t-1}$ & $z_{t-2}$ & $z_{t-3}$ & & $z_{t-29}$\\
\midrule
\vHMatUma \vHMatUmb \vHMatUmc \vHMatUmd
\multicolumn{8}{c}{$\vdots$ \quad (\vHRealObsUm{} linhas)}\\
\bottomrule
\end{tabular}
\end{center}
Repare na diagonal: o $z_t$ de uma linha reaparece como $z_{t-1}$ na linha de baixo. É a mesma série, deslocada.

\subsection{O que sai dele}
O MQO dá os $\vHRealM$ coeficientes do AR longo. Com eles, o resíduo de cada dia,
\[
  \hat\varepsilon_t = z_t - \bigl(\hat a_0 + \hat a_1 z_{t-1} + \dots + \hat a_m z_{t-m}\bigr),
\]
é a \textbf{estimativa do erro} daquele dia. Isso é tudo o que se quer do estágio 1: os coeficientes do AR
longo são jogados fora.

Na série simulada dá para conferir, porque ali o erro verdadeiro é conhecido (Figura~\ref{fig:erros}). A
correlação entre o erro estimado e o verdadeiro foi de $\vHSintCorrUm$.

\begin{figure}[htbp]\centering
\begin{tikzpicture}
\begin{axis}[base, name=t, width=9.6cm, height=5.6cm, xmin=179, xmax=238, ylabel={Erro},
  xlabel={Dia da série simulada}, title={Os últimos 60 dias},
  legend style={at={(0.5,-0.32)}, anchor=north, font=\footnotesize, legend columns=2}]
\addplot[cdad, line width=1.4pt] table[x=t, y=eps]{\arq{sint_serie.csv}};
\addlegendentry{erro verdadeiro $\varepsilon_t$}
\addplot[cari, line width=0.9pt, dashed] table[x=t, y=eps_est]{\arq{sint_serie.csv}};
\addlegendentry{estimado no estágio 1 $\hat\varepsilon_t$}
\end{axis}
\begin{axis}[base, at={(t.east)}, anchor=west, xshift=1.4cm, width=5.8cm, height=5.8cm, xmin=-3.5, xmax=3.5, ymin=-3.5, ymax=3.5,
  xlabel={$\varepsilon_t$ verdadeiro}, ylabel={$\hat\varepsilon_t$ estimado}, title={Todos os dias\\correlação $\vHSintCorrUm$}]
\draw[black!40, dashed] (axis cs:-3.5,-3.5) -- (axis cs:3.5,3.5);
\addplot[only marks, mark=*, mark size=1pt, cmod, opacity=0.6] table[x=eps, y=eps_est]{\arq{sint_serie.csv}};
\end{axis}
\end{tikzpicture}
\caption{Série simulada: os erros estimados pelo AR longo seguem os verdadeiros de perto. Não são iguais, e
é por isso que o resultado do Hannan-Rissanen é uma aproximação.}
\label{fig:erros}
\end{figure}

%=====================================================================
\section{Estágio 2: a regressão que dá os coeficientes}

Agora a coluna que faltava existe. O estágio 2 regride $z_t$ contra a constante, os retornos passados e os
\emph{erros estimados} passados:
\[
  z_t = c + \phi_1 z_{t-1} + \dots + \phi_p z_{t-p} + \theta_1 \hat\varepsilon_{t-1} + \dots + \theta_q \hat\varepsilon_{t-q} + \text{resto}.
\]
É de novo um MQO comum, com $1 + p + q$ colunas. Os primeiros $m + q$ retornos ficam de fora (no código,
\texttt{primeiro = max(m + q, p)}), porque antes deles não há $\hat\varepsilon$ suficientes. No pregão real, com $q = 1$,
ficam de fora os primeiros $\vHRealPrimeiro$ retornos e a regressão tem $\vHRealObsDois$ linhas:

\begin{center}\small
\begin{tabular}{l r | c r}
\toprule
dia $t$ & $z_t$ & 1 & $\hat\varepsilon_{t-1}$\\
\midrule
\vHMatDoisa \vHMatDoisb \vHMatDoisc \vHMatDoisd
\multicolumn{4}{c}{$\vdots$ \quad (\vHRealObsDois{} linhas)}\\
\bottomrule
\end{tabular}
\end{center}

Os coeficientes dessa regressão são o resultado: $c$, os $\phi$ e os $\theta$ do modelo. Na série simulada deram
$\hat\phi = \vHSintPhiEst$ e $\hat\theta = \vHSintThetaEst$, contra $\vHSintPhi$ e $\vHSintTheta$ verdadeiros. É uma
única amostra: a seção~\ref{sec:montecarlo} mostra o comportamento em mil.

\begin{ideia}[{O fluxo completo, como está no \texttt{CArima::HannanRissanen}}]
\begin{center}
\begin{tikzpicture}[seta/.style={-{Stealth[length=5pt]}, line width=0.8pt, black!60}]
\node[fluxo=cdad] (z) {log-retornos\\$z_1, \dots, z_N$};
\node[fluxo=cmod, right=0.55cm of z] (e1) {\textbf{estágio 1}\\AR($m$) por MQO\\$\Rightarrow$ $\hat\varepsilon_t$};
\node[fluxo=cmod, right=0.55cm of e1] (e2) {\textbf{estágio 2}\\$z_t$ contra $z_{t-i}$ e $\hat\varepsilon_{t-j}$\\$\Rightarrow$ $c, \phi, \theta$};
\node[fluxo=cok, right=0.55cm of e2] (r) {\textbf{resíduos finais}\\recalculados com\\$c, \phi, \theta$};
\draw[seta] (z) -- (e1); \draw[seta] (e1) -- (e2); \draw[seta] (e2) -- (r);
\node[fluxo=cdad, below=0.6cm of e1] (atalho) {sem MA ($q = 0$):\\o estágio 1 é pulado};
\draw[seta, dashed] (z.south) |- (atalho.west);
\draw[seta, dashed] (atalho.east) -| (e2.south);
\node[font=\footnotesize, align=left, anchor=west, text width=3.4cm] at ([xshift=0.4cm]atalho.east -| e2.east) {os resíduos finais vão\\para a previsão, o AIC\\e o GARCH};
\end{tikzpicture}
\end{center}
\end{ideia}

%=====================================================================
\section{Depois dos dois estágios: os resíduos finais}

Um detalhe que costuma passar despercebido. Os $\hat\varepsilon_t$ do estágio 1 \emph{não} são os resíduos que o robô usa
depois. Com os coeficientes finais em mãos, o código (\texttt{CArima::CalcularResiduos}) recalcula todos os erros
do início, um dia de cada vez:
\[
  e_t = z_t - \bigl(c + \phi_1 z_{t-1} + \dots + \theta_1 e_{t-1} + \dots\bigr),
\]
começando com $e = 0$ nos primeiros dias. São esses $e_t$ que vão para três lugares: a previsão do dia seguinte,
a verossimilhança que o AIC compara na escolha das ordens e o GARCH, que estima a volatilidade em cima deles.

Na série simulada, os resíduos finais tiveram correlação de $\vHSintCorrFinal$ com o erro verdadeiro, mais do que os
do estágio 1 ($\vHSintCorrUm$). Com o modelo certo e poucos coeficientes, a conta recursiva acerta mais que o AR
longo, que carrega o ruído de $\vHSintM$ coeficientes.

%=====================================================================
\section{Um pregão real do seu robô: \vHDiaReal}\label{sec:real}

Este é o pregão de \vHDiaReal{} da configuração ARIMA-GARCH Sharpe, no traço da refação de 05/10/2026. Nesse dia o
modelo em uso era um \textbf{ARIMA(0,1,1) com constante}: nenhum retorno passado, um erro passado. É o caso mais
simples em que os dois estágios rodam, e por isso o estágio 2 cabe num gráfico.

\subsection{Passo a passo}
\begin{enumerate}[leftmargin=*, itemsep=2pt]
\item \textbf{Dados.} Os 240 fechamentos antes do pregão, de \vHJanIni{} a \vHJanFim{} (último: R\$ \vHRealFechUlt).
  O log de cada um, e a diferença: $N = \vHRealN$ log-retornos.
\item \textbf{Estágio 1.} AR($\vHRealM$) com constante, por MQO, em $\vHRealObsUm$ linhas. O desvio-padrão dos resíduos
  foi de $\vHRealDpUm\%$ ao dia. A Figura~\ref{fig:real}(c) mostra os $\vHRealM$ coeficientes.
\item \textbf{Estágio 2.} $z_t = c + \theta\,\hat\varepsilon_{t-1}$, em $\vHRealObsDois$ linhas. Com uma coluna só além da
  constante, é uma reta num gráfico de dispersão (Figura~\ref{fig:real}b): $c = \vHRealC\%$ e $\theta = \vHRealTheta$.
\item \textbf{Resíduos finais.} Recalculados com $c$ e $\theta$. O do último dia da janela foi $e_N = \vHRealEN\%$.
\item \textbf{Previsão.}
  \[
    \hat z_{N+1} = c + \theta\, e_N = \vHRealC\% + (\vHRealTheta)\times(\vHRealEN\%) = \vHRealC\% + \vHRealParteMa\% = \vHRealPrev\%.
  \]
  Positiva, então o ARIMA deu sinal de compra.
\end{enumerate}

\begin{figure}[htbp]\centering
\begin{tikzpicture}
\begin{groupplot}[base, group style={group size=2 by 2, horizontal sep=3.1cm, vertical sep=2.9cm}]
\nextgroupplot[width=7.4cm, height=5.6cm, ylabel={\% ao dia}, xlabel={Dia da janela},
  title={(a) Os últimos 60 retornos e os erros}, xmin=179, xmax=238, ymin=-10, ymax=10, clip=true,
  legend style={at={(0.5,-0.3)}, anchor=north, font=\footnotesize, legend columns=3}]
\addplot[cdad, line width=1.3pt] table[x=i, y=z]{\arq{real_serie.csv}};
\addlegendentry{retorno $z_t$}
\addplot[cari, dashed] table[x=i, y=eps_est]{\arq{real_serie.csv}};
\addlegendentry{$\hat\varepsilon_t$ (estágio 1)}
\addplot[cmod] table[x=i, y=eps_final]{\arq{real_serie.csv}};
\addlegendentry{$e_t$ (final)}
\nextgroupplot[width=6cm, height=6cm, xmin=-12, xmax=12, ymin=-12, ymax=12,
  xlabel={$\hat\varepsilon_{t-1}$, erro estimado de ontem (\%)}, ylabel={$z_t$, retorno de hoje (\%)},
  title={(b) O estágio 2 inteiro}]
\addplot[only marks, mark=*, mark size=1.1pt, cmod, opacity=0.6] table[x=eps_ant, y=z]{\arq{real_estagio2.csv}};
\addplot[cdes, line width=1.3pt, domain=-12:12, samples=2] {\pHRealC + \pHRealTheta*x};
\node[cdes, font=\footnotesize, anchor=north east, fill=white, fill opacity=0.85, text opacity=1] at (rel axis cs:0.98,0.98)
  {inclinação $\theta = \vHRealTheta$};
\nextgroupplot[width=7.4cm, height=5cm, ybar, bar width=3.5pt, xmin=0, xmax=30, xlabel={Defasagem $j$},
  ylabel={Coeficiente}, title={(c) O AR($\vHRealM$) do estágio 1}, /pgf/number format/fixed]
\addplot[fill=cari!80, draw=none] table[x=j, y=coef]{\arq{real_pesos.csv}};
\nextgroupplot[width=6cm, height=5cm, xbar, bar width=9pt, xmin=-0.05, xmax=0.3,
  ytick={1,2,3}, yticklabels={constante $c$, parte MA $\theta e_N$, previsão}, ymin=0.4, ymax=3.6,
  xlabel={\% ao dia}, title={(d) De onde vem a previsão}, /pgf/number format/fixed, /pgf/number format/precision=3,
  nodes near coords, nodes near coords style={font=\scriptsize}]
\addplot[fill=cmod!70, draw=none] coordinates {(\pHRealC,1) (\pHRealParteMa,2) (\pHRealPrev,3)};
\end{groupplot}
\end{tikzpicture}
\caption{O pregão de \vHDiaReal, refeito com as mesmas regras do robô. Em (b), cada ponto é um dia: o erro
estimado do dia anterior contra o retorno do dia. A nuvem é quase redonda e a reta, quase deitada: o erro de
ontem diz pouco sobre o retorno de hoje. Em (a), um dia de queda forte sai cortado na escala.}
\label{fig:real}
\end{figure}

\subsection{O que esse dia ensina}
O $\theta$ é pequeno e a nuvem da Figura~\ref{fig:real}(b) é quase redonda. Ainda assim, nesse pregão a parte MA
pesou \emph{mais} que a constante na previsão ($\vHRealParteMa\%$ contra $\vHRealC\%$), porque o erro de ontem foi grande
($\vHRealEN\%$) e o $\theta$ negativo transforma um erro negativo em empurrão para cima. É a leitura de
``reversão'': depois de um dia pior que o esperado, o modelo espera uma pequena compensação.

A máxima verossimilhança, na mesma janela, daria $\theta = \vHRealMvTheta$ e previsão de $\vHRealMvPrev\%$: números
diferentes, mesmo sinal, mesma decisão.

\begin{ideia}[Conferência]
Os valores de $c$, $\theta$ e da previsão refeitos aqui são \textbf{idênticos}, em todas as casas gravadas, aos que o
robô registrou no traço desse pregão (\texttt{refeito\_2026-10-05/traco\_demo/traco\_demo\_GARCH.csv}). O mesmo vale
para o pregão de 05/06/2020, com um ARIMA(0,1,4) e quatro $\theta$. O relatório descreve o que o robô de fato fez.
\end{ideia}

%=====================================================================
\section{Onde o Hannan-Rissanen trabalha no robô}

O método entra em dois momentos, com pesos bem diferentes.

\textbf{Todo pregão}, na estimação do modelo em uso. Aqui os dois estágios só rodam se o modelo tiver termo MA
($q > 0$). Com $q = 0$ o código pula o estágio 1 e faz um MQO direto, e com $p = q = 0$ a constante é só a média
dos retornos da janela. No traço de jul/2019 a jun/2021, os dois estágios rodaram em \vHTracoDois{} pregões e o atalho
em \vHTracoUm{} (Figura~\ref{fig:traco}).

\textbf{A cada 21 pregões}, na escolha das ordens. A grade testa os 25 modelos com $p$ e $q$ de 0 a 4. Os 20 que têm
$q \ge 1$ passam pelos dois estágios, e o AIC de cada um é calculado com os resíduos finais da seção anterior. É
aqui que o método mais trabalha: ele decide, indiretamente, se o robô vai usar um modelo com MA ou não.

\begin{figure}[htbp]\centering
\begin{tikzpicture}
\begin{axis}[base, date coordinates in=x, width=14.5cm, height=5.4cm, xmin=2019-07-01, xmax=2021-06-30,
  xtick={2019-07-01,2020-01-01,2020-07-01,2021-01-01,2021-06-30}, xticklabels={jul/2019,jan/2020,jul/2020,jan/2021,jun/2021},
  ymin=-0.3, ymax=0.3, ylabel={$\theta_1$ em uso}, /pgf/number format/fixed,
  title={O $\theta_1$ estimado a cada pregão, só nos dias com termo MA}]
\addplot[black!40, line width=0.5pt] coordinates {(2019-07-01,0) (2021-06-30,0)};
\addplot[only marks, mark=*, mark size=1.2pt, cmod] table[x=data, y expr={\thisrow{q}>0 ? \thisrow{theta1} : nan}]{\arqr{dia_a_dia.csv}};
\end{axis}
\end{tikzpicture}
\caption{Traço da configuração ARIMA-GARCH Sharpe. Cada ponto é um pregão em que o modelo em uso tinha termo
MA e os dois estágios rodaram. Os blocos correspondem aos períodos entre duas escolhas de ordem. Nos dias sem
ponto o modelo não tinha MA.}
\label{fig:traco}
\end{figure}

%=====================================================================
\section{Quão bem ele estima: mil simulações}\label{sec:montecarlo}

Uma série só não diz se o método é bom. Para isso, gerou-se mil séries com coeficientes conhecidos, cada uma do
tamanho da janela do robô ($\vHSintN$ retornos), e estimou-se cada uma de três jeitos: Hannan-Rissanen (o do robô),
máxima verossimilhança (a opção experimental \texttt{MetodoARIMA=1}, que parte do Hannan-Rissanen e refina) e um
AR(1) que ignora o MA.

\begin{figure}[htbp]\centering
\begin{tikzpicture}
\begin{groupplot}[base, group style={group size=3 by 1, horizontal sep=1.3cm, ylabels at=edge left},
  width=5.6cm, height=5cm, ymin=0, ylabel={Réplicas}, /pgf/number format/fixed]
\nextgroupplot[title={$\hat\phi$ no ARMA(1,1)\\(verdadeiro $\vHSintPhi$)}, xmin=0.3, xmax=0.95,
  legend style={at={(0.02,0.98)}, anchor=north west, font=\scriptsize}]
\addplot[hist={bins=30, data min=0.3, data max=0.95}, fill=cmod!50, draw=cmod, line width=0.3pt] table[y=hr_phi]{\arq{mc_arma.csv}};
\addlegendentry{HR}
\addplot[hist={bins=30, data min=0.3, data max=0.95}, fill=cari!40, draw=cari, opacity=0.7, line width=0.3pt] table[y=ingenuo]{\arq{mc_arma.csv}};
\addlegendentry{AR(1) sem MA}
\draw[cdes, line width=1.2pt] (axis cs:\pHSintPhi,0) -- (axis cs:\pHSintPhi,300);
\nextgroupplot[title={$\hat\theta$ no ARMA(1,1)\\(verdadeiro $\vHSintTheta$)}, xmin=0, xmax=0.8,
  legend style={at={(0.98,0.98)}, anchor=north east, font=\scriptsize}]
\addplot[hist={bins=30, data min=0, data max=0.8}, fill=cmod!50, draw=cmod, line width=0.3pt] table[y=hr_theta]{\arq{mc_arma.csv}};
\addlegendentry{HR}
\addplot[hist={bins=30, data min=0, data max=0.8}, draw=cok, line width=1pt, fill=none] table[y=mv_theta]{\arq{mc_arma.csv}};
\addlegendentry{máx. verossim.}
\draw[cdes, line width=1.2pt] (axis cs:\pHSintTheta,0) -- (axis cs:\pHSintTheta,300);
\nextgroupplot[title={$\hat\theta$ num MA(1) com o $\theta$\\do pregão real ($\vHRealTheta$)}, xmin=-0.4, xmax=0.2]
\addplot[hist={bins=30, data min=-0.4, data max=0.2}, fill=cmod!50, draw=cmod, line width=0.3pt] table[y=hr_theta]{\arq{mc_ma.csv}};
\draw[cdes, line width=1.2pt] (axis cs:\pHRealTheta,0) -- (axis cs:\pHRealTheta,300);
\draw[black, dashed, line width=0.8pt] (axis cs:0,0) -- (axis cs:0,300);
\end{groupplot}
\end{tikzpicture}
\caption{Distribuição das estimativas em \vHMcArmaReplicas{} séries simuladas. A linha vermelha é o valor verdadeiro;
a tracejada preta, no gráfico da direita, é o zero.}
\label{fig:mc}
\end{figure}

\begin{center}\small
\begin{tabular}{l c c c}
\toprule
& \multicolumn{2}{c}{ARMA(1,1), $\phi = \vHSintPhi$, $\theta = \vHSintTheta$} & MA(1), $\theta = \vHRealTheta$\\
\cmidrule(lr){2-3}\cmidrule(lr){4-4}
Método & $\hat\phi$ média (desvio) & $\hat\theta$ média (desvio) & $\hat\theta$ média (desvio)\\
\midrule
Hannan-Rissanen (robô) & $\vHMcArmaHrPhi$ ($\vHMcArmaHrPhiDp$) & $\vHMcArmaHrTheta$ ($\vHMcArmaHrThetaDp$) & $\vHMcMaHrTheta$ ($\vHMcMaHrThetaDp$)\\
Máxima verossimilhança & $\vHMcArmaMvPhi$ ($\vHMcArmaMvPhiDp$) & $\vHMcArmaMvTheta$ ($\vHMcArmaMvThetaDp$) & $\vHMcMaMvTheta$ ($\vHMcMaMvThetaDp$)\\
AR(1) ignorando o MA & $\vHMcArmaIngenuo$ ($\vHMcArmaIngenuoDp$) & -- & --\\
\bottomrule
\end{tabular}
\end{center}

\textbf{Ignorar o MA distorce.} O AR(1) sem MA erra o $\phi$ de forma sistemática (média $\vHMcArmaIngenuo$ para um
verdadeiro de $\vHSintPhi$). É a justificativa para ter o termo MA e, portanto, um método que o estime.

\textbf{O HR chega perto, com um pouco mais de erro que a máxima verossimilhança.} No $\theta$, o desvio do HR foi
$\vHMcArmaHrThetaDp$ e o da máxima verossimilhança, $\vHMcArmaMvThetaDp$, e a média do HR ficou um pouco abaixo do
verdadeiro. É o comportamento descrito na literatura: Gómez e Maravall (1998, p.~29) dizem que os estimadores do
segundo estágio são consistentes, mas têm viés e não são os mais eficientes.

\textbf{No tamanho dos coeficientes do robô, o erro de estimação é do tamanho do coeficiente.} Com $\theta = \vHRealTheta$
(o valor que o robô estimou em \vHDiaReal), o desvio da estimativa foi $\vHMcMaHrThetaDp$, quase igual ao próprio
$\theta$. Em $\vHMcMaSinalCerto\%$ das séries o HR acertou o sinal, e nas demais estimou um $\theta$ com o sinal trocado. Isso
não é defeito do método: com 239 retornos, um efeito desse tamanho mal se distingue do zero, e a máxima
verossimilhança tem o mesmo problema ($\vHMcMaMvThetaDp$).

%=====================================================================
\section{O seu código comparado ao método da literatura}

O código segue a ideia de Hannan e Rissanen, mas com algumas escolhas próprias. Comparando com a descrição de
Gómez e Maravall (1998, p.~28--29):

\begin{center}\small
\begin{tabular}{>{\raggedright}p{3.2cm} >{\raggedright}p{5.3cm} >{\raggedright\arraybackslash}p{5.6cm}}
\toprule
Ponto & Gómez e Maravall (1998) & O seu código (\texttt{Arima.mqh})\\
\midrule
Estágio 1 só com MA & sim & sim (com $q = 0$, MQO direto)\\
Como estima o AR longo & autocorrelações e recursão de Durbin-Levinson & MQO com constante\\
Ordem do AR longo & $\max(\lfloor\log^2(N)\rfloor,\ 2\max(p,q))$ & $\max(\lfloor\ln^2 N\rfloor,\ p+q+1)$, no máximo $N/5$\\
Terceiro passo (correção do viés) & sim, quando há MA & \textbf{não}\\
Critério para escolher $(p,q)$ & BIC & AIC sobre amostra comum\\
Invertibilidade ($|\theta|<1$) & -- & não é verificada\\
\bottomrule
\end{tabular}
\end{center}

O que isso significa na prática:
\begin{itemize}[leftmargin=*, itemsep=2pt]
\item \textbf{Sem o terceiro passo}, o robô usa os coeficientes do estágio 2 como estão. É de onde vem o pequeno
  viés do $\theta$ nas simulações. A máxima verossimilhança faria o papel desse refinamento, e foi testada como
  opção em 15/10/2026 e não adotada.
\item \textbf{MQO no lugar de Durbin-Levinson} dá coeficientes parecidos, mas não iguais, para o AR longo. Os dois são
  formas aceitas de estimar um AR.
\item \textbf{A invertibilidade não é imposta.} Nos dias do traço os $\theta$ ficaram pequenos (Figura~\ref{fig:traco}),
  longe de 1, então isso não pesou ali.
\end{itemize}

\subsection{Um ponto para pensar: a ordem do AR longo numa amostra pequena}
A regra $\ln^2 N$ vem de resultados para amostras grandes. Com $N = \vHSintN$ ela dá $\vHSintM$ defasagens, estimadas com
$\vHSintObsUm$ linhas. Na série simulada, refazendo tudo com cada ordem de 1 a $\vHSintMMax$, a qualidade dos erros
estimados caiu à medida que a ordem cresceu (Figura~\ref{fig:ordem}). Com poucas defasagens o AR longo já capta um
ARMA(1,1). Cada defasagem a mais acrescenta um coeficiente estimado com ruído.

\begin{figure}[htbp]\centering
\begin{tikzpicture}
\begin{groupplot}[base, group style={group size=2 by 1, horizontal sep=1.6cm}, width=7.4cm, height=5.4cm,
  xmin=0, xmax=48, xlabel={Ordem $m$ do AR longo}, /pgf/number format/fixed]
\nextgroupplot[ylabel={Correlação com o erro verdadeiro}, title={Qualidade dos erros estimados}, ymin=0.86, ymax=1]
\addplot[cmod] table[x=m, y=corr]{\arq{sint_ordem.csv}};
\draw[cdes, dashed] (axis cs:\pHSintM,0.86) -- (axis cs:\pHSintM,1) node[pos=0.1, right, font=\footnotesize] {regra: $m = \vHSintM$};
\nextgroupplot[ylabel={Estimativa}, title={Coeficientes finais}, ymin=0.3, ymax=0.8,
  legend style={at={(0.98,0.97)}, anchor=north east, font=\footnotesize, legend columns=2}]
\addplot[cmod] table[x=m, y=phi]{\arq{sint_ordem.csv}};
\addlegendentry{$\hat\phi$}
\addplot[cari] table[x=m, y=theta]{\arq{sint_ordem.csv}};
\addlegendentry{$\hat\theta$}
\addplot[cmod, dotted, line width=0.8pt, forget plot] coordinates {(0,\pHSintPhi) (48,\pHSintPhi)};
\addplot[cari, dotted, line width=0.8pt, forget plot] coordinates {(0,\pHSintTheta) (48,\pHSintTheta)};
\draw[cdes, dashed] (axis cs:\pHSintM,0.3) -- (axis cs:\pHSintM,0.8);
\end{groupplot}
\end{tikzpicture}
\caption{Uma única série simulada (ARMA(1,1), $\phi = \vHSintPhi$, $\theta = \vHSintTheta$), refeita com cada ordem do AR longo.
As linhas pontilhadas são os valores verdadeiros. A correlação foi medida sobre os mesmos dias em todas as
ordens.}
\label{fig:ordem}
\end{figure}

\begin{cuidado}[Até onde isso vale]
É uma série só, com coeficientes bem maiores que os da PETR4. Não se testou se uma ordem menor mudaria alguma
decisão do robô. Fica registrado como ponto em aberto (\texttt{ia/INCERTEZAS.md}, item 84), não como defeito.
\end{cuidado}

%=====================================================================
\section{Perguntas que a banca pode fazer}

\begin{banca}
\textbf{Por que Hannan-Rissanen e não máxima verossimilhança?}\\
Porque o modelo é reestimado todo pregão e, na escolha das ordens, 25 vezes de uma vez. O Hannan-Rissanen são
duas regressões lineares: não precisa de ponto de partida, de taxa de aprendizado nem de número de iterações,
não depende da escala dos preços e não tem problema de convergência. A literatura o usa justamente na
identificação automática de modelos (Gómez e Maravall, 1998, p.~10). O custo é uma estimativa um pouco menos
precisa: nas simulações, o desvio do $\theta$ foi $\vHMcArmaHrThetaDp$, contra $\vHMcArmaMvThetaDp$ da máxima verossimilhança.

\medskip
\textbf{O que é o estágio 1, em uma frase?}\\
Um modelo autorregressivo longo, ajustado por mínimos quadrados, cujos resíduos servem de estimativa dos erros que
o termo de médias móveis precisa e que não podem ser observados.

\medskip
\textbf{Por que um AR longo consegue estimar os erros de um modelo com MA?}\\
Porque todo ARMA invertível pode ser escrito como um AR de ordem infinita, com pesos que diminuem. Um AR longo,
mas finito, aproxima bem esse AR infinito.

\medskip
\textbf{O ``MA'' do ARIMA é a média móvel da análise técnica?}\\
Não. O MA do ARIMA é uma combinação dos \emph{erros} passados, não dos preços. A ponte com a média móvel de Almeida
(2016) é pela parte AR, não pela parte MA.

\medskip
\textbf{Como vocês sabem que a implementação está certa?}\\
Duas conferências. Em série simulada com coeficientes conhecidos, o método recupera os valores verdadeiros dentro
do erro esperado. E a estimação de pregões reais, refeita fora do robô com os mesmos passos, dá os mesmos
coeficientes e a mesma previsão que o robô registrou.
\end{banca}

%=====================================================================
\section*{Glossário rápido}
\begin{description}[leftmargin=2.6cm, style=nextline, itemsep=1pt, font=\normalfont\bfseries]
\item[MQO] Mínimos quadrados ordinários: escolhe os coeficientes de uma regressão que tornam mínima a soma dos
  erros ao quadrado. Conta fechada, sem iteração.
\item[AR($p$)] Modelo autorregressivo: o valor de hoje depende dos $p$ valores anteriores.
\item[MA($q$)] Médias móveis (de erros): o valor de hoje depende dos $q$ erros anteriores.
\item[Invertível] Um MA cujos erros podem ser reescritos como soma de valores passados ($|\theta| < 1$ num termo).
\item[Resíduo] A diferença entre o valor observado e o que o modelo ajustado previa. É a estimativa do erro.
\item[Viés] Tendência de uma estimativa a errar sempre para o mesmo lado.
\item[Monte Carlo] Repetir um experimento muitas vezes com dados simulados, para ver como um método se comporta
  em média.
\end{description}

\vfill
\noindent\rule{\linewidth}{0.4pt}
{\footnotesize\textbf{Procedência.} Dados gerados por \nolinkurl{ferramentas/HannanRissanenDetalhado.mq5} (rodado por
\nolinkurl{ferramentas/scripts/hannan_rissanen.sh}) em \nolinkurl{dados/resultados/hannan_rissanen_2026-10-30/}, e o
traço da refação de 05/10/2026 resumido em \nolinkurl{dados/resultados/robo_em_acao_2026-10-30/}. Semente das
simulações: 20261030. Obras citadas: Gómez e Maravall (1998), em \nolinkurl{bibliografia/pdfs/Gomez-1998.pdf};
Hannan e Rissanen (1982) é citado através deles (o artigo original não foi lido). Material de estudo, não faz parte
do TG.}
\end{document}
